%% kapa: exam=1 round=24 year=2013 period=1 form=A subject=civil_law subject_ko=민법 seq=25 number=25
%% source: Q-Net 감정평가사 기출문제 (cst00309) artlSeq=5073802 file=extracted/1008/5073802/2013년 제24회 감정평가사 1차시험 기출문제/1교시(A형).hwp
%% conversion: kordoc 4.24.51 -> tools/gen1.py (auto); review=proofread-sonnet; flags=none
\begin{question}{25}
선의취득에 관한 설명으로 옳지 않은 것은? (다툼이 있으면 판례에 의함)
\begin{choices}
  \item 등기ㆍ등록으로 공시되는 자동차, 선박의 선의취득은 인정되지 않는다.
  \item 점유개정에 의한 선의취득은 부정되지만 목적물반환청구권의 양도에 의한 선의취득은 가능하다.
  \item 점유보조자가 횡령한 물건은 민법상 도품에 해당하지 않는다.
  \item 선의취득을 위해서는 양수인의 점유취득이 선의․무과실로 이루어져야 하는데, 선의․무과실의 기준시점은 물권적 합의와 인도 중에서 어느 것이라도 먼저 갖추어진 때이다.
  \item 선의취득의 요건을 갖춘 경우 양수인은 선의취득의 효과를 부정하고 임의로 종전소유자에게 동산을 받아갈 것을 요구할 수 없다.
\end{choices}
\end{question}
\answer{4}
